\scanpage{92}
% L1, underlined
\uline{Fonction modulaire}

% L2–10
En considérant $g_2,g_3$ comme fonctions de $\omega_1,\omega_2$, où $\unsure{\mathfrak I|J}\frac{\omega_2}{\omega_1}>0$, qui sont homogènes de degré $2$ et $3$ respectivement, on peut former la fonction $\frac{g_2^3}{g_2^3-27g_3^2}$ (finie puisque $\Delta=g_2^3-27g_3^2\ne0$), qui est homogène de degré zéro, donc ne dépend que de $\tau=\frac{\omega_2}{\omega_1}$. On \unsure{l'appelle} fonction modulaire, notation $J(\tau)$; $J(\tau)$ est holomorphe pour $\mathfrak I\tau>0$, et par définition, on a

% L11
\[
\text{pour }\tau=\frac{\omega_2}{\omega_1}\qquad
J(\tau)=\unsure{{}|-}\frac{g_2^3}{\Delta}=\frac1{1-27\frac{g_3^2}{g_2^3}}
\]

% L12–21
Comme $g_2$ et $g_3$ peuvent être pris arbitraires, à \deleted{\unsure{fait|part}} la condition $\Delta=g_2^3-27g_3^2\ne0$ près, $J(\tau)$ prend \deleted{\unsure{toutes}} les valeurs; de plus, le fait que $g_2$ et $g_3$ ne dépendent en fait que du groupe engendré par $\omega_1$ et $\omega_2$ s'exprime par le fait que $J(\tau)$ est invariant par \unsure{les|le} \deleted{\unsure{subst|sous}} groupe arithmétique unimodulaire -- groupe des substitutions $\tau\to\frac{a\unsure{\tau|z}+b}{c\tau+d}$, où $a,b,c,d$ sont entiers, et $ad-bc=1$. De plus,

% lower left
% archive sheet number: 91
